\documentclass[12pt,a4paper]{article}
\usepackage[utf8]{inputenc}
\usepackage[indonesian]{babel}
\usepackage{amsmath}
\usepackage{amsfonts}
\usepackage{amssymb}
\usepackage{geometry}

\geometry{margin=2.5cm}

\title{Penyelesaian Titik Berat Busur Astroida}
\author{}
\date{}

\begin{document}

\maketitle

\section{Soal}
Tentukan titik berat busur astroida yang berada di kuadran I dengan persamaan:
\begin{equation}
  x^{2/3} + y^{2/3} = a^{2/3}
\end{equation}

\section{Penyelesaian}

\subsection{Turunan Pertama ($\frac{dy}{dx}$)}
Turunkan persamaan secara implisit terhadap $x$:
\begin{align*}
  \frac{d}{dx}(x^{2/3} + y^{2/3}) &= \frac{d}{dx}(a^{2/3}) \\
  \frac{2}{3}x^{-1/3} + \frac{2}{3}y^{-1/3} \frac{dy}{dx} &= 0 \\
  \frac{dy}{dx} &= -\frac{x^{-1/3}}{y^{-1/3}} = -\left( \frac{y}{x} \right)^{1/3}
\end{align*}

\subsection{Elemen Panjang Busur ($ds$)}
Gunakan rumus elemen panjang busur $ds = \sqrt{1 + \left(\frac{dy}{dx}\right)^2} dx$:
\begin{align*}
  ds &= \sqrt{1 + \left( -\frac{y^{1/3}}{x^{1/3}} \right)^2} dx \\
  ds &= \sqrt{1 + \frac{y^{2/3}}{x^{2/3}}} dx = \sqrt{\frac{x^{2/3} + y^{2/3}}{x^{2/3}}} dx
\end{align*}
Substitusikan $x^{2/3} + y^{2/3} = a^{2/3}$:
\begin{equation}
  ds = \sqrt{\frac{a^{2/3}}{x^{2/3}}} dx = \frac{a^{1/3}}{x^{1/3}} dx = a^{1/3} x^{-1/3} dx
\end{equation}

\subsection{Panjang Busur Total ($L$)}
Integralkan $ds$ dari $x = 0$ sampai $x = a$:
\begin{align*}
  L &= \int_{0}^{a} a^{1/3} x^{-1/3} dx \\
  L &= a^{1/3} \left[ \frac{3}{2} x^{2/3} \right]_0^a = a^{1/3} \left( \frac{3}{2} a^{2/3} \right) = \frac{3}{2}a
\end{align*}

\subsection{Koordinat Titik Berat ($\bar{x}$)}
\begin{align*}
  \bar{x} &= \frac{1}{L} \int x \, ds \\
  \bar{x} &= \frac{1}{\frac{3}{2}a} \int_{0}^{a} x (a^{1/3} x^{-1/3}) dx \\
  \bar{x} &= \frac{2 a^{1/3}}{3a} \int_{0}^{a} x^{2/3} dx \\
  \bar{x} &= \frac{2 a^{1/3}}{3a} \left[ \frac{3}{5} x^{5/3} \right]_0^a \\
  \bar{x} &= \frac{2 a^{1/3}}{3a} \left( \frac{3}{5} a^{5/3} \right) = \frac{2}{5} \frac{a^2}{a} = \frac{2}{5}a
\end{align*}

Berdasarkan sifat simetri kurva terhadap garis $y = x$ di kuadran I, maka $\bar{y} = \bar{x}$.

\subsection{Koordinat Titik Berat ($\bar{y}$)}
Untuk membuktikan secara eksplisit tanpa asumsi simetri:
\begin{equation*}
  \bar{y} = \frac{1}{L} \int y \, ds
\end{equation*}
Substitusikan $y = (a^{2/3} - x^{2/3})^{3/2}$ dan $ds = a^{1/3} x^{-1/3} dx$:
\begin{equation*}
  \bar{y} = \frac{2 a^{1/3}}{3a} \int_{0}^{a} (a^{2/3} - x^{2/3})^{3/2} x^{-1/3} dx
\end{equation*}
Gunakan substitusi $u = a^{2/3} - x^{2/3}$, sehingga $du = -\frac{2}{3}x^{-1/3} dx$:
\begin{align*}
  \bar{y} &= \frac{2 a^{1/3}}{3a} \int_{a^{2/3}}^{0} u^{3/2} \left( -\frac{3}{2} du \right) \\
  \bar{y} &= \frac{a^{1/3}}{a} \int_{0}^{a^{2/3}} u^{3/2} du \\
  \bar{y} &= a^{-2/3} \left[ \frac{2}{5} u^{5/2} \right]_0^{a^{2/3}} \\
  \bar{y} &= a^{-2/3} \left( \frac{2}{5} a^{5/3} \right) = \frac{2}{5} a
\end{align*}

\section{Kesimpulan}
Titik berat busur astroida di kuadran I adalah $(\bar{x}, \bar{y}) = \mathbf{\left(\frac{2}{5}a, \frac{2}{5}a\right)}$.

\end{document}